\section*{Chapter \ref{ch:b2:investment} --- Investment}

\begin{solution}{exo:b2:investment:1}
% figdata: grunfeld_summary
That investment rises with $q$, here with market value relative to the firm's average: the fitted slope is 0.93 and the correlation
0.57. The wide spread at a given value, and the two thirds of the variation the correlation leaves, are not explained:
firm-specific shocks, the gap between average and marginal $q$, and perhaps finance.
\end{solution}

\begin{solution}{exo:b2:investment:2}
% code: exercises/solutions.py user_cost
$200 \times (0.04 + 0.06 - 0.01) = 200 \times 0.09 = 18$ a year.
\end{solution}

\begin{solution}{exo:b2:investment:3}
% figdata: iy_summary
In 1973, at 36.5 per cent. Before 1914 the averages are 16.9 (United States), 12.8 (Germany) and 15.5 per cent (Japan); since 1950, 20.3,
23.4 and 27.7: all higher, Japan's most.
\end{solution}

\begin{solution}{exo:b2:investment:4}
A variable known to the firm, such as cash flow or sales, with a clear effect on investment once $q$ is held fixed, in a
large panel and with $q$ well measured. The debate box reports such a cash-flow effect, but disputes whether $q$ is well
measured.
\end{solution}

\begin{solution}{exo:b2:investment:5}
% code: exercises/solutions.py desired_capital
The stock falls. The user cost becomes $0.07 + 0.10 = 0.17$, and $K^* = 0.3 \times 1{,}000/0.17 = 1{,}764.7$, about 12 per cent less.
\end{solution}

\begin{solution}{exo:b2:investment:6}
% code: exercises/solutions.py q_investment
$0.08 + 0.3/10 = 0.11$ and $0.08 - 0.1/10 = 0.07$. With $q < 1$ installed capital is worth less than its replacement cost: the firm invests less than
it loses to depreciation, and its capital shrinks.
\end{solution}

\begin{solution}{exo:b2:investment:7}
% code: exercises/solutions.py grunfeld
Pooled: 0.115 on value and 0.228 on capital, $R^2 = 0.82$. With a constant for each firm: 0.110 (standard error 0.011) and 0.310 (0.017),
$R^2 = 0.95$. The firm constants absorb permanent differences between firms, such as their size and the kind of capital they
use, so the second regression uses only changes within each firm, the variation $q$ theory speaks about.
\end{solution}

\begin{solution}{exo:b2:investment:8}
Today $q$ jumps up, by less than for an immediate rise, and investment starts at once. Until the rise takes effect, the
old $\dot q = 0$ curve still governs: $q$ keeps rising and capital grows, and the path reaches the new saddle path exactly when the
rise takes effect. Then $q$ and capital move along it to the new steady state, $q$ falling back to 1.
\end{solution}

\begin{solution}{exo:b2:investment:9}
% code: exercises/solutions.py now, wait
Now: $11 \times 10 - 100 = 10$. Waiting: in the good state profit is 12.5 and the project is worth $137.5 - 100 = 37.5$; in the bad state, 7.5 and $82.5 < 100$, so no
investment. Waiting is worth $\tfrac12 \times 37.5/1.1 = 17.05 > 10$: wait. Consistently, 10 is below the trigger 11.48.
\end{solution}

\begin{solution}{exo:b2:investment:10}
One more unit of investment costs $1 + \chi(I/K - \delta)$, the derivative of $I + \tfrac{\chi}{2}(I/K - \delta)^2K$ with respect to $I$, and adds one unit of installed capital, worth $q$. The firm
invests until the two are equal: $1 + \chi(I/K - \delta) = q$, so $I/K = \delta + (q - 1)/\chi$.
\end{solution}

\begin{solution}{exo:b2:investment:11}
% code: exercises/solutions.py trigger
Building now is worth $Lx - C$; waiting is worth $(Lx(1 + u) - C)/k$, since only the good state builds, with probability one half, a year later. At the
trigger $Lx - C = (Lx(1 + u) - C)/k$, so $Lx(1 - (1 + u)/k) = C(1 - 1/k)$, which gives the formula. As $u$ rises, $1 - (1 + u)/k$ falls (it stays positive while $u < 1 + 2r$), so $x^*$
rises: with $C = 100$ and $r = 0.10$, 9.09, 11.48, 15.58 and 24.24 for $u = 0$, 0.25, 0.5, 0.75.
\end{solution}

\begin{solution}{exo:b2:investment:12}
Investment follows marginal $q$. If measured $q$ is marginal $q$ plus noise, the regression's coefficient on measured $q$ is
too small, and any variable correlated with true $q$ picks up part of its effect. Cash flow today that \omterm{def:b2:signalling-and-screening:signal}{signals} future
profits is such a variable: it predicts investment given measured $q$, for a firm that could borrow all it wanted.
\end{solution}

\begin{solution}{pb:b2:investment:1}
% code: exercises/solutions.py terminal
\textbf{Part I.}
1. $50 \times (0.06 + 0.08 - 0.02) = 6$ a year.
2. $20 - K = 6$: 14 cranes.
3. $50 \times (0.08 + 0.08 - 0.02) = 7$; $20 - K = 7$: 13 cranes.
4. $50 \times (0.06 + 0.08) = 7$.
5. $45 \times 0.12 = 5.4$; $K = 14.6$, so 14 whole cranes: the fifteenth earns 5, less than 5.4.
6. A crane loses 8 per cent of its value each year: the terminal bears that loss whether or not it pays anything out.
\textbf{Part II.}
7. $900/(10 \times 50) = 1.8$.
8. Constant returns to scale in production and in \omterm{def:b2:investment:adjustment-cost}{adjustment costs}, and prices the terminal takes as given.
9. $0.08 + 0.8/10 = 0.16$: 1.6 cranes a year.
10. $0.08$: 0.8 crane a year, replacing wear.
11. $0.08 - 0.2/10 = 0.06$: 0.6 crane, less than wear.
12. \omterm{def:b2:investment:adjustment-cost}{Adjustment costs} rise faster than the speed of investment: buying four at once costs far more than buying them over
several years.
\textbf{Part III.}
13. $11 \times 12 - 100 = 32$.
14. $\tfrac12 \times (11 \times 18 - 100)/1.1 = 98/2.2 = 44.55$; in the bad state $11 \times 6 = 66 < 100$, so no berth.
15. $44.55 - 32 = 12.55$.
16. Wait: waiting is worth more.
17. Build now: if the bad state comes it sells the berth at its cost, so waiting protects against nothing.
18. $\tfrac12 \times (11 \times 15 - 100)/1.1 = 65/2.2 = 29.55 < 32$: build now.
\textbf{Part IV.}
19. $100/11 = 9.09$.
20. 15.58 (the listing).
21. 24.24.
22. Each raises its trigger, so fewer cross it, and those that do wait for a higher profit: investment falls and responds
less to demand.
23. Investment that does not fall, or whose response to demand does not shrink, when measured uncertainty rises, among
firms whose investment is hard to reverse.
24. At a profit of 12 and $u = 0.5$, the \omterm{def:b2:investment:option-value}{option value of waiting} is 12.55: more than a third of the net present value of 32. The terminal
builds at once only when this year's profit reaches 15.58, 1.7 times the 9.09 at which the net present value turns positive.
\end{solution}
